% ============================================================
\chapter{Group Theory for Deep Learning}
\label{ch:groups}
% ============================================================

\section{Introduction to Group Theory}

Symmetry is everywhere. A snowflake looks the same after a $60°$ rotation. The laws of physics do not change if you translate your laboratory ten metres to the left. A molecule's chemical properties depend on which rearrangements of its atoms leave its structure invariant. But what \emph{is} a symmetry, precisely? The answer was forged in the early nineteenth century by two brilliant and tragically short-lived mathematicians: \'Evariste Galois, who died in a duel at twenty, and Niels Henrik Abel, who succumbed to tuberculosis at twenty-six. Their legacy is \emph{group theory}---the mathematical language of symmetry.

A group is nothing more than a set of transformations that can be composed, reversed, and that include ``doing nothing.'' Yet from this minimal definition springs an astonishingly rich theory. For geometric deep learning, groups are the key to building neural networks that \emph{respect} the symmetries of their input data: if rotating an image should not change its classification, the network's architecture should encode rotational symmetry from the start. This chapter builds the algebraic foundations that make such constructions possible.

\begin{definition}[Group]
A \textbf{group} is a set $G$ equipped with a binary operation
$\cdot: G \times G \to G$ satisfying:
\begin{enumerate}
    \item \textbf{Associativity}: $(g \cdot h) \cdot k = g \cdot (h \cdot k)$
          for all $g, h, k \in G$.
    \item \textbf{Identity}: there exists $e \in G$ such that
          $e \cdot g = g \cdot e = g$ for all $g \in G$.
    \item \textbf{Inverse}: for every $g \in G$, there exists $g^{-1} \in G$ such
          that $g \cdot g^{-1} = g^{-1} \cdot g = e$.
\end{enumerate}
\end{definition}

\begin{example}[Fundamental groups in GDL]
\begin{enumerate}
    \item \textbf{Symmetric group} $S_n$: permutations of $n$ elements.
          $\abs{S_n} = n!$.
    \item \textbf{Cyclic group} $\Z_n = \Z/n\Z$: discrete rotations.
    \item \textbf{Orthogonal group} $\mathrm{O}(n)$: matrices $\mathbf{Q} \in \R^{n \times n}$
          with $\mathbf{Q}^\top\mathbf{Q} = \mathbf{I}$.
    \item \textbf{Special orthogonal group} $\mathrm{SO}(n)$: rotations,
          $\det(\mathbf{Q}) = +1$.
    \item \textbf{Euclidean group} $\mathrm{E}(n) = \mathrm{O}(n) \ltimes \R^n$:
          rotations, reflections, and translations.
    \item \textbf{Special Euclidean group} $\mathrm{SE}(n) = \mathrm{SO}(n) \ltimes \R^n$:
          rotations and translations (rigid motions).
\end{enumerate}
\end{example}

\subsection{Subgroups and Homomorphisms}

\begin{definition}[Subgroup]
A subset $H \subseteq G$ is a \textbf{subgroup} of $G$ (denoted $H \leq G$) if
$H$ is itself a group under the operation of $G$.
\end{definition}

\begin{definition}[Group homomorphism]
A \textbf{homomorphism} of groups is a map $\varphi: G \to H$ such that
$\varphi(g_1 \cdot g_2) = \varphi(g_1) \cdot \varphi(g_2)$ for all
$g_1, g_2 \in G$. If $\varphi$ is bijective, it is an \textbf{isomorphism}.
\end{definition}

\begin{proposition}[Properties of homomorphisms]
Let $\varphi: G \to H$ be a homomorphism. Then:
\begin{enumerate}
    \item $\varphi(e_G) = e_H$.
    \item $\varphi(g^{-1}) = \varphi(g)^{-1}$ for all $g \in G$.
    \item $\ker(\varphi) = \{g \in G : \varphi(g) = e_H\}$ is a normal subgroup of $G$.
\end{enumerate}
\end{proposition}

\section{Group Actions}

\begin{definition}[Group action]
A (left) \textbf{action} of a group $G$ on a set $X$ is a map
$\cdot: G \times X \to X$ such that:
\begin{enumerate}
    \item $e \cdot x = x$ for all $x \in X$.
    \item $(g \cdot h) \cdot x = g \cdot (h \cdot x)$ for all $g, h \in G$ and $x \in X$.
\end{enumerate}
\end{definition}

\begin{example}[Actions in deep learning]
\begin{enumerate}
    \item $S_n$ acts on $\R^n$ by permuting coordinates:
          $\pi \cdot (x_1, \ldots, x_n) = (x_{\pi^{-1}(1)}, \ldots, x_{\pi^{-1}(n)})$.
    \item $\mathrm{SO}(3)$ acts on $\R^3$ by rotation: $R \cdot \mathbf{x} = R\mathbf{x}$.
    \item $\Z^2$ acts on images by translation: $\mathbf{t} \cdot f(\mathbf{x}) = f(\mathbf{x} - \mathbf{t})$.
\end{enumerate}
\end{example}

\begin{definition}[Orbit and stabilizer]
Let $G$ act on $X$. For $x \in X$:
\begin{itemize}
    \item The \textbf{orbit} of $x$ is $G \cdot x = \{g \cdot x : g \in G\}$.
    \item The \textbf{stabilizer} of $x$ is $G_x = \{g \in G : g \cdot x = x\}$.
\end{itemize}
\end{definition}

\begin{theorem}[Orbit--stabilizer theorem]
For a finite group $G$ acting on $X$ and $x \in X$:
\[
    \abs{G} = \abs{G \cdot x} \cdot \abs{G_x}
\]
\end{theorem}

\section{Group Representations}

\begin{definition}[Representation]
A \textbf{representation} of a group $G$ on a vector space $V$ is a
homomorphism $\rho: G \to \mathrm{GL}(V)$, i.e., a map assigning to each
$g \in G$ an invertible linear transformation $\rho(g): V \to V$ such that
$\rho(g_1 g_2) = \rho(g_1)\rho(g_2)$.
\end{definition}

\begin{example}[Common representations]
\begin{enumerate}
    \item \textbf{Trivial representation}: $\rho(g) = \mathrm{Id}$ for all $g$.
    \item \textbf{Permutation representation}: $\rho(\pi) = \mathbf{P}_\pi$
          (permutation matrix).
    \item \textbf{Standard representation of $\mathrm{SO}(3)$}: $\rho(R) = R$
          (the rotation matrix itself).
    \item \textbf{Regular representation}: $[\rho(g)f](h) = f(g^{-1}h)$.
\end{enumerate}
\end{example}

\subsection{Irreducible Representations}

\begin{definition}[Irreducible representation]
A representation $\rho: G \to \mathrm{GL}(V)$ is \textbf{irreducible} if the
only subspaces of $V$ invariant under the action of $G$ are $\{0\}$ and $V$.
\end{definition}

\begin{theorem}[Maschke's theorem]
Every finite-dimensional representation of a finite group (over $\R$ or $\C$)
can be decomposed as a direct sum of irreducible representations:
\[
    V = V_1 \oplus V_2 \oplus \cdots \oplus V_k
\]
where each $V_i$ is an irreducible subspace.
\end{theorem}

\begin{keyformulas}[title=\textbf{Irreducible representations of $\mathrm{SO}(3)$}]
The irreducible representations of $\mathrm{SO}(3)$ are indexed by
$\ell \in \{0, 1, 2, \ldots\}$ and have dimension $2\ell + 1$:
\begin{itemize}
    \item $\ell = 0$: scalars (dimension 1)---trivial representation.
    \item $\ell = 1$: vectors (dimension 3)---standard representation.
    \item $\ell = 2$: traceless symmetric tensors (dimension 5).
\end{itemize}
The representation matrices are the \textbf{Wigner D-matrices}
$D^\ell_{mm'}(R)$, related to spherical harmonics:
\[
    Y_\ell^m(R^{-1}\hat{\mathbf{r}}) = \sum_{m'=-\ell}^{\ell}
    D^\ell_{mm'}(R)\, Y_\ell^{m'}(\hat{\mathbf{r}})
\]
\end{keyformulas}

\section{Equivariance and Invariance}

\begin{definition}[Equivariant map]
Let $\rho_1: G \to \mathrm{GL}(V_1)$ and $\rho_2: G \to \mathrm{GL}(V_2)$
be two representations. A linear map $T: V_1 \to V_2$ is \textbf{equivariant}
(or an \textbf{intertwiner}) if:
\[
    T \circ \rho_1(g) = \rho_2(g) \circ T, \quad \forall g \in G.
\]
When $\rho_2$ is the trivial representation, $T$ is \textbf{invariant}.
\end{definition}

\begin{theorem}[Schur's lemma]
Let $\rho_1$ and $\rho_2$ be two irreducible representations of $G$ and
$T: V_1 \to V_2$ an intertwiner:
\begin{enumerate}
    \item If $\rho_1 \not\cong \rho_2$, then $T = 0$.
    \item If $\rho_1 \cong \rho_2$ (over $\C$), then $T = \lambda\, \mathrm{Id}$
          for some $\lambda \in \C$.
\end{enumerate}
\end{theorem}

\begin{intuition}[title=\textbf{Importance of Schur's lemma}]
Schur's lemma is fundamental because it \textbf{strongly constrains} the form
of equivariant layers in a neural network. If equivariance is imposed, the
space of allowed linear maps is drastically reduced, which justifies weight
sharing.
\end{intuition}

\section{Tensor Products of Representations}

\begin{definition}[Tensor product]
The \textbf{tensor product} of two representations $\rho_1: G \to \mathrm{GL}(V_1)$
and $\rho_2: G \to \mathrm{GL}(V_2)$ is the representation
$\rho_1 \otimes \rho_2: G \to \mathrm{GL}(V_1 \otimes V_2)$ defined by:
\[
    (\rho_1 \otimes \rho_2)(g)(v_1 \otimes v_2) = \rho_1(g)v_1 \otimes \rho_2(g)v_2.
\]
\end{definition}

\begin{theorem}[Clebsch--Gordan decomposition for $\mathrm{SO}(3)$]
The tensor product of two irreducible representations of $\mathrm{SO}(3)$
decomposes as:
\[
    D^{\ell_1} \otimes D^{\ell_2} = \bigoplus_{\ell=\abs{\ell_1-\ell_2}}^{\ell_1+\ell_2} D^{\ell}
\]
The change-of-basis coefficients are the \textbf{Clebsch--Gordan coefficients}
$C^{\ell m}_{\ell_1 m_1, \ell_2 m_2}$.
\end{theorem}

\begin{codeblock}[title=\textbf{Tensor product with e3nn}]
\begin{minted}{python}
import torch
from e3nn import o3

# Irreducible representations of SO(3)
irreps_1 = o3.Irreps("1x0e + 1x1o")  # scalar + vector
irreps_2 = o3.Irreps("1x1o")          # vector

# Tensor product
tp = o3.FullTensorProduct(irreps_1, irreps_2)
print(f"Input 1: {irreps_1}")
print(f"Input 2: {irreps_2}")
print(f"Output: {tp.irreps_out}")
# 0e x 1o -> 1o, 1o x 1o -> 0e + 1e + 2e

x1 = irreps_1.randn(5, -1)  # batch of 5
x2 = irreps_2.randn(5, -1)
out = tp(x1, x2)
print(f"Output shape: {out.shape}")
\end{minted}
\end{codeblock}

\section{Lie Groups}

\begin{definition}[Lie group]
A \textbf{Lie group} is a group $G$ that is also a smooth manifold, such that
the multiplication $(g, h) \mapsto gh$ and inversion $g \mapsto g^{-1}$
operations are smooth.
\end{definition}

\begin{example}[Lie groups in GDL]
\begin{center}
\begin{tabular}{llll}
\toprule
\textbf{Group} & \textbf{Dimension} & \textbf{Compact?} & \textbf{Usage} \\
\midrule
$\mathrm{SO}(2)$ & 1 & Yes & 2D rotations \\
$\mathrm{SO}(3)$ & 3 & Yes & 3D rotations \\
$\mathrm{SE}(3)$ & 6 & No & Rigid motions \\
$\mathrm{O}(3)$ & 3 & Yes & Rotations + reflections \\
$\mathrm{GL}(n, \R)$ & $n^2$ & No & Linear transformations \\
\bottomrule
\end{tabular}
\end{center}
\end{example}

\subsection{Lie Algebra}

\begin{definition}[Lie algebra]
The \textbf{Lie algebra} $\mathfrak{g}$ of a Lie group $G$ is the tangent space
to $G$ at the identity $e$, equipped with the \textbf{Lie bracket}
$[\cdot, \cdot]: \mathfrak{g} \times \mathfrak{g} \to \mathfrak{g}$.
\end{definition}

\begin{keyformulas}[title=\textbf{Exponential map}]
The \textbf{exponential map} $\exp: \mathfrak{g} \to G$ connects the Lie
algebra to the group:
\[
    \exp(X) = \sum_{k=0}^{\infty} \frac{X^k}{k!}
\]
For $\mathrm{SO}(3)$, the Lie algebra $\mathfrak{so}(3)$ is the space of
$3 \times 3$ skew-symmetric matrices. Rodrigues' formula gives:
\[
    \exp(\theta\, [\hat{\mathbf{n}}]_\times) = \mathbf{I} + \sin\theta\,
    [\hat{\mathbf{n}}]_\times + (1 - \cos\theta)\, [\hat{\mathbf{n}}]_\times^2
\]
where $[\hat{\mathbf{n}}]_\times$ is the skew-symmetric matrix of the unit
vector $\hat{\mathbf{n}}$.
\end{keyformulas}

\section{Group Convolution}

\begin{definition}[Group convolution]
Let $G$ be a locally compact group with Haar measure $\mu$. The
\textbf{group convolution} of two functions $f, g: G \to \R$ is:
\[
    (f * g)(x) = \int_G f(y)\, g(y^{-1}x)\, d\mu(y).
\]
\end{definition}

\begin{proposition}[Equivariance of convolution]
Group convolution is equivariant to left translation: if
$L_a f(x) = f(a^{-1}x)$, then:
\[
    L_a(f * g) = (L_a f) * g.
\]
This property justifies the use of convolution as the fundamental layer in
equivariant networks.
\end{proposition}

\begin{theorem}[CNN as convolution on $\Z^2$]
A classical convolutional network performs a group convolution on
$G = \Z^2$ (the discrete translation group). The first layer performs a
cross-correlation:
\[
    [f \star \psi](x) = \sum_{y \in \Z^2} f(y)\, \psi(y - x)
    = \sum_{y \in \Z^2} f(x + y)\, \psi(y)
\]
which is equivariant to translations of $f$.
\end{theorem}

\section{Group Equivariant CNNs (G-CNNs)}

Cohen and Welling (2016) generalized CNNs by replacing the translation
group $\Z^2$ with a larger group $G$.

\begin{definition}[G-CNN]
A \textbf{G-CNN} is a network whose layers perform convolutions on a
group $G$. The first layer \textit{lifts} the signal from $\Z^2$ to $G$:
\[
    [\text{lift}(f)](g) = \sum_{x \in \Z^2} f(x)\, \psi(g^{-1} \cdot x)
\]
Subsequent layers operate entirely on $G$:
\[
    [f * \psi](g) = \sum_{h \in G} f(h)\, \psi(h^{-1}g)
\]
\end{definition}

\begin{codeblock}[title=\textbf{G-CNN for group $p4$ (90\textdegree{} rotations)}]
\begin{minted}{python}
import torch
import torch.nn as nn

class P4Conv(nn.Module):
    """Equivariant convolution for group p4 (0/90/180/270 rotations)."""
    def __init__(self, in_channels, out_channels, kernel_size=3):
        super().__init__()
        self.weight = nn.Parameter(
            torch.randn(out_channels, in_channels, kernel_size, kernel_size)
        )
        nn.init.kaiming_normal_(self.weight)

    def _rotate_kernel(self, w, k):
        """Rotate kernel by k*90 degrees."""
        return torch.rot90(w, k, dims=[-2, -1])

    def forward(self, x):
        # x: (batch, C_in, H, W) or (batch, C_in, 4, H, W)
        outputs = []
        for k in range(4):  # 4 rotations
            w_rot = self._rotate_kernel(self.weight, k)
            out = torch.nn.functional.conv2d(
                x if x.dim() == 4 else x.sum(dim=2),
                w_rot, padding=1
            )
            outputs.append(out)
        # Stack on group dimension
        return torch.stack(outputs, dim=2)  # (batch, C_out, 4, H, W)

# Test
conv = P4Conv(3, 16)
x = torch.randn(2, 3, 32, 32)
out = conv(x)
print(f"Output: {out.shape}")  # (2, 16, 4, 32, 32)
\end{minted}
\end{codeblock}

\section{Linear Invariants and Equivariants}

\begin{theorem}[Characterization of equivariant linear maps]
The space of linear maps $T: V_1 \to V_2$ that are equivariant with respect
to representations $\rho_1$ and $\rho_2$ is:
\[
    \mathrm{Hom}_G(V_1, V_2) = \{T \in \mathrm{Hom}(V_1, V_2) :
    T\rho_1(g) = \rho_2(g)T,\; \forall g \in G\}
\]
Its dimension is given by:
\[
    \dim \mathrm{Hom}_G(V_1, V_2) = \inner{\chi_1}{\chi_2}_G
    = \frac{1}{\abs{G}} \sum_{g \in G} \overline{\chi_1(g)}\, \chi_2(g)
\]
where $\chi_i(g) = \tr(\rho_i(g))$ is the \textbf{character} of $\rho_i$.
\end{theorem}

\begin{corollary}[Number of free parameters]
For an equivariant layer between two representations decomposed into
irreducibles $V_1 = \bigoplus_i n_i V_i^{\text{irr}}$ and
$V_2 = \bigoplus_j m_j V_j^{\text{irr}}$, the number of free parameters is:
\[
    \sum_k n_k \cdot m_k
\]
where the sum is over common irreducible types.
\end{corollary}

\section*{Exercises}

\begin{exercise}[Symmetry groups of molecules]
\begin{enumerate}
    \item Determine the symmetry group of the water molecule $\mathrm{H_2O}$
          (group $C_{2v}$).
    \item Enumerate all irreducible representations of $C_{2v}$.
    \item Explain how these symmetries constrain a molecular energy prediction model.
\end{enumerate}
\end{exercise}

\begin{exercise}[Representations of $S_3$]
\begin{enumerate}
    \item Enumerate the six elements of the symmetric group $S_3$.
    \item Find all irreducible representations of $S_3$ (there are three).
    \item Verify the relation $\sum_i (\dim V_i)^2 = \abs{G}$.
    \item Decompose the standard permutation representation $\R^3$ into irreducibles.
\end{enumerate}
\end{exercise}

\begin{exercise}[Rodrigues' formula]
\begin{enumerate}
    \item Verify Rodrigues' formula for $\theta = \pi/2$ and
          $\hat{\mathbf{n}} = (0, 0, 1)$.
    \item Show that $\exp(\theta [\hat{\mathbf{n}}]_\times) \in \mathrm{SO}(3)$.
    \item Implement the exponential map $\mathfrak{so}(3) \to \mathrm{SO}(3)$ in
          PyTorch and verify orthogonality of the result.
\end{enumerate}
\end{exercise}

\begin{exercise}[Discrete group convolution]
Implement a group convolution on the dihedral group $D_4$ (symmetries of
the square).
\begin{enumerate}
    \item Enumerate the 8 elements of $D_4$.
    \item Implement the multiplication table of $D_4$.
    \item Implement the group convolution and verify its equivariance.
\end{enumerate}
\end{exercise}
